\chapter*{Abstract}
\addcontentsline{toc}{chapter}{Abstract}

\section*{Abstract}
The reliability of avionics systems is critical to aviation safety and operational efficiency. Unexpected failures lead to costly disruptions and safety risks. This project presents an intelligent software platform for real-time health monitoring and predictive maintenance of avionics systems. The platform analyzes time-series data from key parameters including voltage levels, current measurements, temperature sensors, and error rate counters.

Using a combination of statistical methods (threshold-based detection) and machine learning algorithms (Isolation Forest), the system identifies three types of anomalies: sudden spikes, gradual drifts, and intermittent failure patterns. An explainability module provides interpretable alerts by analyzing feature contributions, supporting maintenance decision-making.

The system follows a modular architecture comprising data ingestion, preprocessing, anomaly detection, REST API, and an interactive web dashboard. The project is developed using the Scrum agile methodology, with four sprints delivering incremental functionality from data pipeline to complete prototype.

This prototype demonstrates the feasibility of integrating data-driven anomaly detection into avionics maintenance workflows, paving the way for future industrial deployment with real aircraft data.

\textbf{Keywords:} Avionics, Predictive Maintenance, Anomaly Detection, Machine Learning, Isolation Forest, Time-Series Analysis, Scrum, Python, FastAPI

\bigskip
\section*{Résumé}
La fiabilité des systèmes avioniques est essentielle pour la sécurité aérienne et l'efficacité opérationnelle. Les pannes imprévues entraînent des perturbations coûteuses et des risques de sécurité. Ce projet présente une plateforme logicielle intelligente pour la surveillance de l'état de santé en temps réel et la maintenance prédictive des systèmes avioniques. La plateforme analyse les données temporelles de paramètres clés tels que les niveaux de tension, les mesures de courant, les capteurs de température et les compteurs de taux d'erreur.

En combinant des méthodes statistiques (détection par seuil) et des algorithmes d'apprentissage automatique (Isolation Forest), le système identifie trois types d'anomalies : les pics soudains, les dérives progressives et les pannes intermittentes. Un module d'explicabilité fournit des alertes interprétables en analysant les contributions des caractéristiques, facilitant ainsi la prise de décision en matière de maintenance.

Le système suit une architecture modulaire comprenant l'ingestion de données, le prétraitement, la détection d'anomalies, une API REST et un tableau de bord web interactif. Le projet est développé selon la méthodologie agile Scrum, avec quatre sprints délivrant des fonctionnalités incrémentales allant du pipeline de données au prototype complet.

Ce prototype démontre la faisabilité de l'intégration de la détection d'anomalies basée sur les données dans les flux de travail de maintenance avionique, ouvrant la voie à un futur déploiement industriel avec des données réelles d'aéronefs.

\textbf{Mots-clés :} Avionique, Maintenance Prédictive, Détection d'Anomalies, Apprentissage Automatique, Isolation Forest, Analyse de Séries Temporelles, Scrum, Python, FastAPI
